\clearpage
\section{Proposition 12.2 in every finite environment}
\label{sec:proposition-complete}
\termdef{completed-proposition12-2}{Completed Proposition 12.2}.
The public theorem now proves, for every finite distinct-list variable
environment $\Gamma$ and atom environment $\Delta$,
\[
 \Gamma\mid\Delta\vDash
 \mathtt{let}\ x=\mathtt{ref}\ 1\ \mathtt{in}\
 \lambda f.\,(f\,();!x)
 \ \simeq^{\ddagger}\ \lambda f.\,(f\,();1).
\]
The Lean statement uses the original scoped terms and their unchanged
Figure~16 translation. Its execution is the complete copying UAM, and
its observations range over every original binding-free closing graph
context at every common adequate fragment allowance. No determinism,
robustness, translation-correctness or callback-termination premise is supplied.

\subsection{The original operational prefix at every allowance}
An operational pair survives one extra allowance by its already proved
padding law. Induction gives a pair after any $n$ extra allowances.
For a pair used inside a fixed certified graph context, first prove the
fragment equivalence at every padded allowance. The existing context-padding
and context-congruence theorems then lift it through that actual context.
The supplied endpoint isomorphisms are padded too. This proves the helper
for the entire contextual use; it does not promote an arbitrary equivalence
to a padded equivalence.
\lean{Copying/PropositionPrefixPadding}{11}{19}{operational_pair_pad}
\lean{Copying/PropositionPrefixPadding}{20}{30}{context_equiv_pad}

Apply this helper to reference creation in the application context containing
the original let-body closure. Apply it again to micro-beta in the certified
store-suffix context. The beta argument remains the stable instance alone,
with its allocated store outside. The intervening store factorization is a
padded recursive isomorphism. Their composition proves the complete
allocation/beta prefix for every $n$. For any $N\geq3$, write $N=3+n$ and
use the proved agreement between adding $n$ allowances and padding to $N$.
\lean{Copying/PropositionPrefixPadding}{31}{39}{original_allocated_pad}
\lean{Copying/PropositionPrefixPadding}{40}{47}{betaSource_opened_pad}
\lean{Copying/PropositionPrefixPadding}{48}{55}{original_opened_pad}
\lean{Copying/PropositionPrefixPadding}{56}{60}{original_opened_padTo}

\subsection{Assembling the closed chain without a padding assumption}
The normalized private graph also commutes with an extra allowance:
serial padding transports its large closure, complete term tree and actual
private cell. Inducting on $d$ gives a recursive isomorphism from the
minimum-depth private closure padded to $d+3$ to $Q_{d+1}$.
Pad the source-normalization calculation itself to that allowance and
append this isomorphism. Exact contextual soundness gives the opened-to-private
comparison at every depth. Compose it with the padded operational prefix,
forgetting costs only where termination equivalence is required.
\lean{Copying/PropositionComplete}{13}{17}{private_raise}
\lean{Copying/PropositionComplete}{19}{29}{private_padding}
\lean{Copying/PropositionComplete}{30}{35}{opened_private_exact}
\lean{Copying/PropositionComplete}{36}{42}{original_private}

\subsection{Arbitrary environments and the final theorem}
The original-environment record is precisely the translated left term padded
to $d+3$. Pad its already proved finite-environment weakening calculation,
and use the proved equality commuting ambient weakening with padding.
This yields exact contextual equivalence to the ambient extension of the
closed original left graph. Lift the complete closed prefix through that
certified ambient context. Finally compose with the completed parametricity
and right-normalization segment in the same environment. The result is
\[
 (P_L^{\Gamma,\Delta})^{\mathrm{pad}(d+3)}
 \simeq
 (P_R^{\Gamma,\Delta})^{\mathrm{pad}(d+3)}.
\]
The cost relation here is universal; the source structural links separately
retain exact successful counts and the earlier operational/parametricity
lemmas retain their directed quantitative comparisons.
\lean{Copying/PropositionComplete}{43}{47}{originalEnvironment}
\lean{Copying/PropositionComplete}{48}{57}{source_environment_exact}
\lean{Copying/PropositionComplete}{58}{75}{equivalence_at}

For arbitrary common allowance $N$, the left term's depth condition gives
$3\leq N$. Writing $N=d+3$ supplies the preceding theorem; the right term's
minimum depth is two. This proves scoped observational equivalence for all
finite scope sizes. The named theorem specializes those sizes to the lengths
of $\Gamma$ and $\Delta$ and accepts their distinctness proofs. The scoped
proof actually works for every pair of natural sizes; names do not enter
the machine. Variables and atoms are separate namespaces, so no disjointness
between the two lists is required. A companion theorem consumes the existing
named-environment record directly.
\lean{Copying/PropositionComplete}{76}{88}{Examples.proposition12_2_scoped}
\lean{Copying/PropositionComplete}{89}{95}{Examples.proposition12_2}
\lean{Copying/PropositionComplete}{96}{100}{Examples.proposition12_2_environment}

\subsection{Semantic regression instances and the scope of the result}
Full concrete term-context congruence lifts the final theorem into any
context of the original annotated grammar. In particular, one regression
passes a callback whose body is the usual self-application loop. The result
compares all observations of the two filled terms without assuming callback
termination. It is not a separate theorem constructing an infinite machine
trace. Another context binds the returned closure once and invokes that same
closure twice with unit callbacks; the same full observational theorem applies.
\lean{Copying/PropositionCompleteChecks}{9}{14}{CompleteChecks.contextual}
\lean{Copying/PropositionCompleteChecks}{15}{16}{CompleteChecks.loopingCallback}
\lean{Copying/PropositionCompleteChecks}{18}{21}{CompleteChecks.loopContext}
\lean{Copying/PropositionCompleteChecks}{22}{26}{CompleteChecks.looping_callback}
\lean{Copying/PropositionCompleteChecks}{27}{29}{CompleteChecks.repeatedContext}
\lean{Copying/PropositionCompleteChecks}{31}{36}{CompleteChecks.repeated_closure_use}

A named mixed-environment instance uses two variable names and three atom
names, deliberately allowing the same spelling in both namespaces. A depth-seven
instance checks an allowance larger than either minimum depth. The final
observation lemma unfolds the theorem to successful-termination equivalence
for every original closing graph context, not only contexts obtained from terms.
\lean{Copying/PropositionCompleteChecks}{37}{42}{CompleteChecks.named_mixed_environment}
\lean{Copying/PropositionCompleteChecks}{43}{48}{CompleteChecks.extra_allowance}
\lean{Copying/PropositionCompleteChecks}{49}{53}{CompleteChecks.observations}

\paragraph{Proof correspondence and remaining audits.}
The proof uses the paper's template method and the indirect copying and
garbage-collection dependencies already discharged by parametricity. It
performs allocation and outer beta before the structural normalization;
Figure~22 displays some structural moves earlier. For sharing components it
uses the proved general sharing-tree theorem, rather than displaying an
elementary associativity/commutativity sequence. All finite representation
maps are kernel-checked, and all semantic changes use proved templates.
The broader correspondence audits retained in Appendix~B concern the
relationship between formal predicates and the paper, and general translation
laws beyond the instances needed here. They are not premises or unfinished
Lean obligations in \texttt{Examples.proposition12\_2}. Thus the formal
endpoint is complete; this does not declare every independent paper-correspondence
question or every theorem in the paper settled.
